\documentclass[conference]{IEEEtran}
\IEEEoverridecommandlockouts
% The preceding line is only needed to identify funding in the first footnote. If that is unneeded, please comment it out.
\usepackage{cite}
\usepackage{amsmath,amssymb,amsfonts}
\usepackage{algorithmic}
\usepackage{graphicx}
\usepackage{textcomp}
\usepackage{xcolor}
\usepackage{float}
\def\BibTeX{{\rm B\kern-.05em{\sc i\kern-.025em b}\kern-.08em
    T\kern-.1667em\lower.7ex\hbox{E}\kern-.125emX}}
\begin{document}

\title{Vn-IPFS: Re-engineering IPFS using von Neumann's Reliable Organism Synthesis Principles}

% \author{\IEEEauthorblockN{Justin Y. Shi}
% \IEEEauthorblockA{\textit{Department of Computer and Information Sciences} \\
% \textit{Temple University}\\
% Philadelphia, PA, USA \\
% shi@temple.edu}
% \and
% \IEEEauthorblockN{Nick Tagliamonte}
% \IEEEauthorblockA{\textit{Department of Computer and Information Sciences} \\
% \textit{Temple University}\\
% Philadelphia, PA, USA \\
% nicholas.tagliamonte@temple.edu}
% \and
% \IEEEauthorblockN{Garrett Stoll}
% \IEEEauthorblockA{\textit{Department of Mathematics} \\
% \textit{Temple University}\\
% Philadelphia, PA, USA \\
% garrett.stoll@temple.edu}
% }

\maketitle

\begin{abstract}
Decentralized storage systems face a persistent tension between global addressability and the overhead of maintaining network state. Traditional Distributed Hash Tables (DHTs) often employ rigid provider-data binding, which limits scalability and data consistency in mission-critical environments. This paper introduces vn-IPFS, a decentralized storage architecture that decouples "provider" and "data" using a probabilistic Tuple Space abstraction. By transforming the storage overlay into an Active Content Addressable Network (ACAN), vn-IPFS implements an N-hop exclusion algorithm for data replication, ensuring survival against regional blackouts through a distribution of near, midrange, and far-flung nodes. The architecture leverages von Neumann’s principles for reliable systems to provide automated failover, partition tolerance via read-repair mechanisms. Experimental results demonstrate that vn-IPFS delivers deterministic performance and high reliability without the scaling constraints of conventional deterministic designs, making it suitable for modern enterprise-grade non-stop storage services.
\end{abstract}

\begin{IEEEkeywords}
decentralized, network, file system, von Neumann, DHT, IPFS
\end{IEEEkeywords}

\section{Introduction}
The landscape of decentralized storage systems has long been defined by the tension between the necessity of global addressability and the fragility of maintaining the state required to achieve it. For the past two decades, the dominant paradigm for solving this problem has been the Distributed Hash Table (DHT) using XOR distances, a class of algorithms that structure the network overlay into a navigable graph \cite{b1}. However, for mission-critical mutable storage systems that require high performance, data consistency, and service reliability, existing DHT systems are not suitable \cite{b2}.

The tight \textless{}provider, data\textgreater{} binding in existing DHT storage systems forces stateful protocol designs that are difficult to optimize. Applying von Neumann principles on reliable organism synthesis using unreliable components \cite{b3}, vn-IPFS proposes to decouple “provider” and “data” while using the same DHT mechanics for efficient decentralized search under Tuple Space semantics. The vn-IPFS architecture is a “probabilistic” architecture without scaling limits, promising to deliver "best effort" non-stop storage services with controlled deterministic performance and reliability.

\begin{itemize}
\item \textbf{The Core Innovation: Provider-Data Decoupling}

The fundamental change in the decentralized storage system design in vn-IPFS is the decoupling of providers from the traditional \textless{}providerID, contentID(CID)\textgreater{} data representation to a \textless{}CID+UID, value\textgreater{} Tuple Space abstraction.
    \begin{itemize}
        \item The Key: Hash of application data content identity combined with user ID.
        \item The Value: A list of providers.
    \end{itemize}
Unlike existing DHT storage systems, the Tuple "put" protocol creates a unique data hash (key) and a unique provider address (IP) hash. The direct data hash key decouples access controls from data store semantics allowing automatic de-duplication (within the owner’s space), data replication with concurrency control, and fault tolerance. This Tuple Space transforms the entire storage overlay network into an Active Content Addressable Network (ACAN).

\item \textbf{Reliability \& Disaster Recovery: Probabilistic vs. Deterministic Systems}

Centralized mission-critical systems suffer fault tolerance and scaling challenges, while traditional decentralized systems are short on sustainable performance. Vn-IPFS bridges this gap, replacing traditional disaster recovery (DR) resources with decentralized equivalents:

\item \textbf{Physical DR Sites vs. Global ACAN}

Traditional DR relies on expensive secondary "hot sites" that can still halt on a single-point failure. In vn-IPFS, there are no idle standby sites; the network can tap into open storage nodes globally under extreme conditions for survival.

\item \textbf{Hardware-Level Replication (RAID) vs. N-Hop Exclusion}

Instead of localized RAID configurations, vn-IPFS replicates each file $\ge 7$ times to randomly selected nodes meeting custom Round Trip Time (RTT) requirements. This uses an application-defined N-hop exclusion algorithm: Near (40\%) for speed, Midrange (30\%) to avoid regional blackouts, and Far-flung (30\%) to sustain global attacks.

\item \textbf{Failover Orchestration vs. Automated Protocol Bypassing}

Failover is handled natively. Node failures are automatically skipped in the protocol. Once a match is found via DHT search, actual data transfer is done through a small direct swarm protocol from multiple providers directly to the requester.

\item \textbf{Split-Brain Management vs. Partition Tolerance}

When a traditional storage network partitions, the system locks up. In vn-IPFS, each partitioned network continues to operate on reachable nodes, allowing updates and deletions. When the network heals, the read protocol repairs missing replicas. For mutable objects, conflict resolution can be customized using cryptographic versioning with timestamps, quorum or any combination of consensus algorithms.
\end{itemize}

Vn-IPFS was designed for fully decentralized high-performance storage without scalability limitations. Routing tables and data replications are completely automated without user interference. In comparison to traditional deterministic architectures, this probabilistic mission-critical architecture delivers deterministic performance with the option to expand into probabilistic behavior under extreme conditions, guaranteeing the "best effort" reliability required by modern enterprises. Preliminary performance results are included.

\begin{table*}[htbp]
\centering
\renewcommand{\arraystretch}{1.5}
\caption{Comparison of Distributed System Architectures}
\label{tab:architecture-comparison}
\begin{tabular}{|c|c|c|c|}
\hline
\textbf{Dimension} & \textbf{Centralized Cloud} & \textbf{\begin{tabular}[c]{@{}c@{}}Legacy DHT\end{tabular}} & \textbf{vn-IPFS} \\ \hline
\textbf{Fault Tolerance} & \begin{tabular}[c]{@{}c@{}}Low \\ (Single-Point)\end{tabular} & High & \begin{tabular}[c]{@{}c@{}}High \\ (No Single-Point Failure)\end{tabular} \\ \hline
\textbf{\begin{tabular}[c]{@{}c@{}}Global \\ Addressability\end{tabular}} & Rigid / Siloed & \begin{tabular}[c]{@{}c@{}}Logarithmic \\ O($\log N$)\end{tabular} & \begin{tabular}[c]{@{}c@{}}Logarithmic O($\log N$) \\ via Tuple Space\end{tabular} \\ \hline
\textbf{\begin{tabular}[c]{@{}c@{}}State \\ Maintenance\end{tabular}} & Heavy / Bottlenecked & \begin{tabular}[c]{@{}c@{}}Stateful Protocol \\ Designs\end{tabular} & \begin{tabular}[c]{@{}c@{}}Stateless Tuple \\ Abstraction\end{tabular} \\ \hline
\textbf{\begin{tabular}[c]{@{}c@{}}Mutation \& \\ Updates\end{tabular}} & Easy & \begin{tabular}[c]{@{}c@{}}Complex / Prone to \\ Race Conditions\end{tabular} & \begin{tabular}[c]{@{}c@{}}Race-Condition-Free \\ via Data Hash Keys\end{tabular} \\ \hline
\textbf{Access Control} & Native / High & Low / Aftermarket & \begin{tabular}[c]{@{}c@{}}Native Centralized \\ Security via \\ Decentralized Processing\end{tabular} \\ \hline
\end{tabular}
\end{table*}

\section{Foundational Technologies }
The fundamental change in the decentralized storage system design in vn-IPFS is the decoupling of providers from the traditional \textless{}providerID, contentID(CID)\textgreater{} data representation to \textless{}CID, value\textgreater{} Tuple Space abstraction for provider-data binding, where the “key” is application data content identity and “value” is the content value(s) hosting a list of providers (or programs when needed). The Tuple Space transforms the entire storage overlay network into an Active Content Addressable Network (or ACAN).

The direct content hash keys decouple the data access and security concerns from storage performance, fault tolerance and deduplication processes. This enables a new dimension in provider-data binding leveraging DHT efficient search and matching machinery to enable race condition-free data replication, load balancing, data consistency controls, and automatic de-duplication. We call this Statistic Multiplexed Communications and Computing (SMC2), which allows the application (in this case, the storage) to exploit all available resources dynamically in real time.

The overall design philosophy is smarter clients and dumb servers with complete end-to-end caretaking at the network edges. This design is a marked departure from the traditional client-server architectures, especially deterministic mission critical architectures.

To incentivize participation in open storage networks, vn-IPFS adopts POS (proof of stake) tokenomics, where the token is named VNFC with 50,000,000,000,000 soft capped supply. Stakers acquire VNFC to become eligible. Once accepted, staked tokens are locked for the duration of the contract until expiration. Stakers do not publicize “storage-available” signals but respond to upload requests when asked. Stakers are selected based on their committed stakes, their RTT from the provider node, storage availability, availability duration, and reputation. The tokenization layer is for incentivizing participation.

The vn-IPFS design allows the storage application to dynamically exploit all resources probabilistically for the best-effort delivery in performance, reliability and security regardless of the network scale.

CAP and Partition tolerance: When the network partitions, many peers are lost. Each partitioned network continues to operate on reachable nodes. Updates and deletions can occur. The read protocol will repair the missing replicas. When the network heals, the peers become live again. For immutable objects (such as tuples), there will be no conflicts, since there is hash uniqueness. For mutable objects, any consensus algorithm can fit the bill, such as cryptographic versioning with timestamp to resolve conflicts. The data owner can also inject custom consensus in real time under extreme conditions.

A mission-critical user would start with a permissioned storage network first.  Expanding the network can enhance system performance and reliability at the same time. The network can also tap into public open storage nodes under extreme conditions for survival.

Like decentralized systems using DHT and variants, there is no single-point failure regardless of how the system scales. The ACAN/SMC2/DHT/p2p protocols ensure better performance, reliability and security as the system expands in size.

The proposed tokenomics encourages stakers sustained support in perpetuity. The token ledger ensures economic benefits for participation. For permissioned networks, they serve as a real time infrastructure monitoring tool. All mission-critical applications will initiate with only permissioned nodes but have the option to tap into open network nodes when needed. Decoupling data access from storage semantics also allows inspection and auditing without replication management complexities. These features make vn-IPFS suitable for all mission-critical webservices for consumer or enterprise applications.

\begin{figure*}[t]
    \centering
    \includegraphics[width=\textwidth]{vnIPFSdecoupling.png}
    \caption{Provider/data decoupling}
    \label{fig:decoupling}
\end{figure*}

\section{Reliability \& Disaster Recovery: Probabilistic vs. Deterministic Systems}

Centralized mission-critical systems suffer fault tolerance and scaling challenges, while traditional decentralized systems are short on sustainable performance. Vn-IPFS bridges this gap, replacing traditional disaster recovery (DR) resources with decentralized equivalents:
\begin{itemize}
    \item Physical DR Sites vs. Global ACAN: Traditional DR relies on expensive secondary "hot sites" that can still halt on a single-point failure. In vn-IPFS, there are no idle standby sites; the network can tap into open storage nodes globally under extreme conditions for survival.

    \item Hardware-Level Replication (RAID) vs. N-Hop Exclusion: Instead of localized RAID configurations, vn-IPFS replicates each file $R \ge 7$ times to randomly selected nodes meeting custom Round Trip Time (RTT) requirements. This uses an application-defined N-hop exclusion algorithm: Near (40\%) for speed, Midrange (30\%) to avoid regional blackouts, and Far-flung (30\%) to sustain global attacks.

    \item Failover Orchestration vs. Automated Protocol Bypassing: Failover is handled natively. Node failures are automatically skipped in the protocol. Once a match is found via DHT search, actual data transfer is done through a small direct swarm protocol from multiple providers directly to the requester.

    \item Split-Brain Management vs. Partition Tolerance: When a traditional network partitions, systems often lock up. In vn-IPFS, each partitioned network continues to operate on reachable nodes, allowing updates and deletions. When the network heals, the read protocol repairs missing replicas. For mutable objects, conflict resolution can be customized using cryptographic versioning with timestamps, quorum or any combination of consensus algorithms.
\end{itemize}

\section{Access Controls}

With data contents encrypted at storage and in transit with owner ID, the proposed storage layer protocol can deliver best effort availability and performance in real time, suitable for mission-critical applications. However, mission-critical systems also require centralized security such as Know Your Customer (KYC) authentication and elaborate access controls that are typically not natively supported by decentralized systems.

The proposed Tuple Space implementation can support centralized security with decentralized processing. The identity layer (accounts) relies on cryptographic keys or wallets natively replicated and synchronized by randomly selected nodes.

Each user account is secured by the public/private key with optional biometric and authorization signatures. Each account can host an entire file system for the user without quota limit. The account is also a crypto wallet holding the user’s current balance. A new user needs to charge the wallet with vnfc coins to initiate the account and the storage system. The storage requests are generated by user’s upload API. Insufficient funds will result in account creation failure or uploading failure. Concurrent account logins are prohibited. This is implemented through the Tuple Space by serliazing requests via the account ID hash. A logged in user can have concurrent access to all folders and files. Concurrent access is also granted for cross-user requests provided each carries the owner’s shared access key.

\section{Put, Read, Delete, Update }

Unlike existing DHT storage systems, the Tuple “put” protocol will create two Tuples: a unique data hash (key) and a unique provider address (IP) hash. The data key decouples access controls from data store semantics allowing automatic de-duplication, data replication and fault tolerance. The value associated with the data key contains the list of provider addresses formed using an application defined N-hop exclusion algorithm to balance access performance and resilience. The application defined N-hop exclusion algorithm should try to meet mission requirements, such as 40\% nearest peers for speed, 30\% mid-range peers to avoid regional blackouts and 30\% far-flung peers to sustain global attacks.  The “put” protocol will try to store the data locally if storage is available or send the data to a local/remote peer when necessary.

The data retrieval “read” and “get/delete” protocols rely on the DHT search engine to identify the matching data globally through XOR distance calculations. Once a match is found, actual data transfer will be done through a small direct swarm protocol from the multiple providers directly to the requester.

Tuple Space keys cannot carry conditions. For mutable storage systems, the direct data hash key is an ideal decoupling device for the storage layer protocols to handle data fault tolerance, data load balancing and owner data de-duplication. Data access controls (security and race/deadlock prevention) are based on data hash keys.

Users either upload or retrieve/update/delete data through the storage system API. Requests are pattern matched via the DHT engine. For a data read request, a small direct data transfer swarm will form between the matching providers and the requestor. Update request is a combination of read+upload. Delete is a combination of read+delete.

\begin{figure*}[t]
    \centering
    \includegraphics[width=\textwidth]{PutReadGet.png}
    \caption{Put, Read, Get Protocols}
    \label{fig:put,read,get}
\end{figure*}

\section{Vn-IPFS Architecture }\label{AA}
\begin{itemize}
    \item P2PTS: A global \textless{}key, value\textgreater{} Tuple Space abstraction implemented using enhanced DHT protocols and supported by SMC2 protocol extension). Node ID = hash(pubkey), each node has a public/private key pair (Ed25519), Message ID (MID) = hash(payload + sender pubkey + nonce). Session key is derived via ECDH. All messages are encrypted with sender pubkey, nonce, and timestamp. Each node maintains a replay protection window. All messages are signed.

    \begin{itemize}
        \item Authenticate Handshake (Prevents MITM \& Identity Forgery):
        \begin{enumerate}
            \item Exchange public keys
            \item Perform challenge-response
            \item Verify signature of random nonce
            \item Perform ECDH
            \item Derive session key
        \end{enumerate}

        \item After handshake:
        \begin{enumerate}
            \item All traffic encrypted (AES-GCM or ChaCha20-Poly1305)
            \item All traffic authenticated
            \item Anti-Replay protection:
            \item Per-peer nonce cache
            \item Reject:
            \begin{itemize}
                \item Duplicate message hash
                \item Expired timestamps
                \item Reused nonce
            \end{itemize}
        \end{enumerate}
    \end{itemize}

    \item API Library:
    \begin{enumerate}
        \item Put(upload). Store in O(1). Replication messages using Bloom filter in O(log N). In this case, false-positives do not matter.
        \item Read. O(lg N) search complexity. Once matched, a small world swarm for quicker load and authentication. Use zk-SNARK non-interactive proofs.
        \item Get(delete).  O(lg N) search complexity with O(1) delete complexity.
    \end{enumerate}

    \item File System parameters:
    \begin{enumerate}
        \item Files are partitioned into 4KB chunks.
        \item All files are encrypted using owner credentials.
        \item Each file is replicated R ($\ge7$) times to randomly selected nodes meeting custom RTT requirements.
        \item All nodes maintain a local routing table: [CID, ProviderID]
    \end{enumerate}

    \item Economic Model:
    \begin{enumerate}
        \item Total Supply Cap: 50,000,000,000 vnfc
        \item Nodes are owned by stakers.
        \item Join the vn-IPFS is free.
        \item Storage contracts are determined by a stake-weighted percentile.
        \item Once the storage contract is concluded, the node simply adds itself to the peer list and informs O(1) random peers.
        \item Token: vnfc
        \item Pricing policy:
        \begin{itemize}
            \item Reputation-weighted awards: uptime, faster retrieval speeds, consistent proof validations and higher reputation scores.
        \end{itemize}
        \item A Percentage of vnfc will be burned decided by DAO consensus.
        \item Tail Emission: Decided by DAO consensus when all reserves are depleted.
    \end{enumerate}

    \item Node Setup and Network Initialization:
    \begin{enumerate}
        \item Hard coded seed nodes in initial peer list as DNS for the system.
        \item All nodes listen on port 5000 (default).
        \item New node joins the network through ver-ack to setup secured communication with the first available peer and expanded peer list to the seed nodes and self. Peer lists are not synchronized amongst all nodes.
        \item Each message carries a unique MID.
        \item Each node maintains a local received MID list (auto expunge every 1-5 minutes), a peer list with IP-addresses, a recent-seen peer list, and a CID routing table (auto updated by “upload” and “delete” protocols).
        \item OpCodes: upload (0:put), read(1:read), delete (2:get) (write = read then upload), inspect(random read)(4), addnode(5)
        \item Replication Factor R: customizable per file/account/group (default $\ge 7$).
        \item Replication distribution is managed in N-hop exclusion protocol: customizable per file/account/group (default: Near(40%), Midrange(30%), Far-flung(30%).
        \item Each node runs a Bitcoin-like peer manager to track abuses: Eclips attacks (limits peers from the same IP range), uses asmap (ASN mapping), maintains diversity in outbound peers; Sybil attacks: bucketing system for address storage, randomized peer eviction logic, DoS attacks: rate limiting, resource use caps and misbehavior scores.
        \item Each node maintains a ban-list of non-trustworthy peers. Ban duration: 24 hours.
    \end{enumerate}

    \item System read protocol on port 5000: [check]
    \begin{enumerate}
        \item Can start on any node but authenticated.
        \item Loop through entire local routing table for near|mid-ranged|far (N|M|F) requirements and actuals.
        \item Send O(lg N) gossip messages with each missing N|M|F requirements entries
        \item If receiving node also finds missing N|M|F entries, do the same. Until none can be found or replay window closes.
    \end{enumerate}

    \item Node join protocol on port 5000: [local IP]
    \begin{enumerate}
        \item Can connect to any node after authentication.
        \item Send O(1) gossip messages to peers.
    \end{enumerate}

    \item Client service on port: 5003
    \begin{enumerate}
        \item Create/update/delete account
        \item Set/update/delete replication vector
        \item Add/read/update/delete wallet
        \item Add/read/update/delete file
        \item App overlay network monitor and refreshing (scripted)
        \item Add/delete node
    \end{enumerate}
\end{itemize}

\section{Performance Evaluation}
Vn-IPFS file upload complexity is O(1) (local storage with amortized replication costs over the propagation paths), but the content key tuple will be logged in O(lg N) steps.

File read complexity is also O(lg N) to find the content key tuple then the provider links to start a download swarm.

File delete complexity is also O(lg N) to find the content key tuple then O(1) for removing the tuple and marking all replicas invalid.

The read protocol will check the integrity of all contents (keys) with associated N|M|F targets in O(lg N). The repair protocol will force updating the local routing tables. This protocol can be triggered anytime by off-line RPC scripts.

Node failures are automatically skipped in protocol. Node “join” protocol has O(1) steps to announce the event and be discovered by the entire network. Multiple nodes can join at arbitrary time in parallel without complications.

Like other DHT-based implementations, the vn-IPFS storage system throughput performance linearly increases as more nodes joining the network incrementally since the increases in protocol complexity of O(lg N) \textless{} O(N).

\begin{figure*}[t]
    \centering
    \includegraphics[width=\textwidth]{vnipfs_graph.png}
    \caption{vn-IPFS Performance}
    \label{fig:vn-IPFS Performance}
\end{figure*}

\begin{figure*}[t]
    \centering
    \includegraphics[width=\textwidth]{swarm_graph.png}
    \caption{SWARM Performance}
    \label{fig:swarmPerformance}
\end{figure*}

\section{Summary}
Vn-IPFS was designed for fully decentralized high-performance storage without scalability limitations. Routing tables and data replications are completely automated without user interference.  Unlike traditional decentralized systems using chartier gossip protocols, the p2pTS maintains efficient routing tables in the entire storage overlay network. Content retrieval complexity is the same O(lg N) as DHT. The customizable replication vector ensures the “best effort” performance and reliability requirements of the user missions. The worst-case retrieval complexity (data not found) is O(lg N) messages. Routing table and content replication and maintenance overheads are amortized over the network and are not directly visible to the users. In comparison to traditional deterministic architectures, the proposed probabilistic mission critical architecture can deliver deterministic performance with the option to expand into probabilistic behavior under extreme conditions.


\begin{thebibliography}{00}
 \bibitem{b1} J. H. Hartman, I. Murdock, and T. Spalink, ``The Swarm scalable storage system,'' in \emph{Proc. 19th IEEE Int. Conf. Distrib. Comput. Syst.}, Austin, TX, USA, 1999, pp. 74-81. doi: 10.1109/ICDCS.1999.776508.

\bibitem{b2} S. Ratnasamy, I. Stoica, and S. Shenker, ``Routing algorithms for DHTs: Some open questions,'' in \emph{Peer-to-Peer Systems (Lecture Notes in Computer Science, vol. 2429)}, P. Druschel, F. Kaashoek, and A. Rowstron, Eds. Berlin, Heidelberg, Germany: Springer, 2002, pp. 45-52. doi: 10.1007/3-540-45748-8\_4.

\bibitem{b3} J. von Neumann, ``Probabilistic logics and the synthesis of reliable organisms from unreliable components,'' Lecture delivered at the California Institute of Technology, Pasadena, CA, USA, Jan. 1952. [Online]. Available: VonNeumann1956.pdf

% \bibitem{b4} Ethereum, ``Ethereum Swarm Sphere.'' swarm.ethereum.org. https://swarm.ethereum.org (accessed May 9, 2026).

% \bibitem{b5} M. Flynn, ``Flynn's taxonomy,'' in \emph{Encyclopedia of Parallel Computing}, Springer, 2011, pp. 689-697. doi: 10.1007/978-0-387-09766-4\_2.

% \bibitem{b6} Wikipedia, ``Fallacies of distributed computing.'' Wikipedia.org. https://en.wikipedia.org/wiki/Fallacies\_of\_distributed\_computing (accessed Oct. 2, 2024).

% \bibitem{b7} G. Amdahl, ``Validity of the single processor approach to achieving large-scale computing capabilities,'' in \emph{Proc. AFIPS Spring Joint Comput. Conf.}, vol. 30, 1967, pp. 483-485. doi: 10.1145/1465482.1465560.

% \bibitem{b8} J. L. Gustafson, ``Reevaluating Amdahl’s law,'' \emph{Commun. ACM}, vol. 32, no. 5, pp. 532-533, 1988. doi: 10.1145/42411.42415.

% \bibitem{b9} Y. Shi, ``Reevaluating Amdahl’s law and Gustafson’s law,'' Temple University, 1996. [Online]. Available: https://cis.temple.edu/\~{}shi/wwwroot/shi/public\_html/docs/amdahl/amdahl.html

% \bibitem{b10} P. Gupta, ``Gossip protocol in distributed systems,'' Medium.com, 2023. [Online]. Available: https://medium.com/nerd-for-tech/gossip-protocol-in-distributed-systems-e2b0665c7135

% \bibitem{b11} D. Loshin, ``Data parallelism,'' in \emph{Business Intelligence}, 2nd ed. ScienceDirect, 2013. [Online]. Available: https://www.sciencedirect.com/topics/computer-science/data-parallelism

% \bibitem{b12} J. Shi, ``Multi-computer system and method,'' U.S. Patent 5 517 656, 1996.

% \bibitem{b13} J. Shi, ``High performance lossless ESB architecture with data protection for mission-critical applications,'' in \emph{Proc. World Cong. Comput. Sci. Inf. Eng.}, vol. 7, 2009.

% \bibitem{b14} J. Shi, ``Statistic multiplexed computing,'' E.P. Patent 3 539 261, 2022.

% \bibitem{b15} J. Shi, ``Statistic multiplexed computing system for network-scale reliable high-performance services,'' U.S. Patent 11 589 926 B1, 2023.

% \bibitem{b16} D. H. Ackley, ``Indefinite scalability for living computation,'' in \emph{Proc. 30th AAAI Conf. Artif. Intell. (AAAI-16)}, 2016.

\end{thebibliography}
\section{Appendix: Peformance Figures}
\end{document}
